\documentclass[12pt]{article}
\usepackage[utf8]{inputenc}
\usepackage[T1]{fontenc}
\usepackage{hyperref}
\title{Zo Language Resources for K-12 Curriculum Development}
\author{Compiled by Claude Sonnet 5.5 (claude-sonnet-5-5) for Andr\'as Kornai}
\date{}
\begin{document}
\maketitle

\section{Introduction}
Zo (ISO 639-3: zom; Glottocode zouu1235) is a Sino-Tibetan (Kuki-Chin) language of Myanmar and India with about 82,000 speakers. The name is easily confused with Tedim (ctd, also called Zomi or Zokam) and with other Chin languages. Searches found very little material that is clearly about Zo itself, and nothing that is a school curriculum.

\section{Primary Educational Resources}
No Zo primer, textbook or school programme turned up in the searches. This section is empty.

\section{Linguistic and Reference Materials}
No Zo-specific dictionary or grammar turned up in the searches. A profile of the Zokam people lists Zo [zom] as related to Tedim Chin and as part of the Sizang branch of the Chin languages. \url{https://www.operation.asiaharvest.org/MYANMAR-Zokam.pdf}

\section{Religious and Community Resources}
The same Asia Harvest profile reports Zo scripture: a New Testament dated 1967, portions dated 1981, a Bible dated 1992, and a Jesus film. The profile is a secondary source, and I did not check the scripture editions themselves.
\url{https://www.operation.asiaharvest.org/MYANMAR-Zokam.pdf}

The Zokam Bible project of the Bible societies concerns the neighbouring Tedim language and should not be treated as a Zo resource without checking: \url{https://www.tbsbibles.org/news/700366/Zokam-Bible-Project.htm}

\section{Digital and Audio Resources}
No Zo-specific audio, corpus or software was found.

\section{Curriculum Development Context}
Zo is reported to be a first language for the whole ethnic community but is not known to be taught in schools. A curriculum would have to start from basic work: identifying the Zo orthography in use, locating the existing scripture text, and building primers. The first step is to check the Zo/Tedim naming overlap, so that Tedim materials are not mistaken for Zo ones.

\end{document}